\chapter{End-to-End Demonstration Workflow}
\label{ch:demo}

This chapter maps a single complete run of \projname{} onto a
recording-friendly workflow suitable for a project demonstration
video (e.g.\ a YouTube walkthrough). Every step below refers to a
capability already implemented in the repository; placeholder figure
frames are marked \emph{[SCREENSHOT]} so that they can be replaced
with real captures without changing the narrative.

\section{Demo Objectives}
\label{sec:demo-obj}

The demonstration should convince a viewer that:

\begin{enumerate}[leftmargin=1.4em]
  \item The platform runs entirely on a laptop with no cloud
        account.
  \item A user can go from raw CSV to reproducible notebook + PDF
        in a single natural-language turn.
  \item The tool trace is fully inspectable.
  \item Every artefact is on the local filesystem and can be opened
        outside the app.
\end{enumerate}

\section{Pre-Recording Checklist}
\label{sec:pre-record}

Before starting the screen recorder:

\begin{itemize}[leftmargin=1.4em]
  \item Confirm \file{.env} contains a valid \code{GROQ\_API\_KEY}.
  \item Empty \file{data/uploads/} and \file{data/artifacts/}
        (leaving only the \file{.gitkeep} sentinels).
  \item Prepare one small classification dataset (Iris or Breast
        Cancer) and one regression dataset (California Housing) in a
        known folder.
  \item Have a browser window sized to 1280$\times$800 for a clean
        capture.
\end{itemize}

\section{Demo Architecture}
\label{sec:demo-arch}

Figure~\ref{fig:demo-arch} shows the deployment topology that the
viewer will see. It is intentionally identical to
Figure~\ref{fig:high-level}; the emphasis for the video is on the
locality of everything except the Groq call.

\begin{figure}[H]
\centering
\begin{tikzpicture}[node distance=0.6cm and 1cm, font=\small\sffamily]
  \node[al-node, minimum width=3.5cm] (br) {Browser};
  \node[al-node, right=of br, minimum width=4cm]
                                                (st) {Streamlit @ 8501};
  \node[al-node, right=of st, minimum width=4cm]
                                                (py) {Python process \\ (pool + tools)};
  \node[al-store, below=of st, minimum width=3.5cm]
                                                (fs) {data/};
  \node[al-node, right=1.4cm of py, fill=alAccent!15, minimum width=3cm]
                                                (llm) {Groq API};

  \draw[al-flow, <->] (br) -- (st);
  \draw[al-flow, <->] (st) -- (py);
  \draw[al-flow, <->] (py.south) |- (fs.east);
  \draw[al-flow, <->] (py) -- (llm);
\end{tikzpicture}
\caption{Deployment for the demonstration recording.}
\label{fig:demo-arch}
\end{figure}

\section{Step-by-Step Walkthrough}
\label{sec:walkthrough}

The recommended script is a \textbf{5-minute} demonstration divided
into eight beats.

\subsection*{Beat 1 (0:00--0:30) --- Launch}
Show a terminal executing
\begin{lstlisting}[style=alBash, numbers=none]
streamlit run app/streamlit_app.py
\end{lstlisting}
and highlight the ``Pool ready (in-process): 5/5 services online, 11
unified tools'' log line. Cut to the browser at
\url{http://localhost:8501}.

\begin{figure}[H]
\centering
\fbox{\begin{minipage}{0.9\linewidth}\centering\vspace{6em}
\emph{[SCREENSHOT 1: Streamlit landing page with sidebar showing all
five services in green.]}\vspace{6em}\end{minipage}}
\caption{Landing screen with service health.}
\label{fig:demo-1}
\end{figure}

\subsection*{Beat 2 (0:30--1:00) --- Upload a dataset}
Drag Iris CSV into the sidebar uploader. Show the success toast and
the ``Active dataset'' selector automatically switching to the new
file.

\begin{figure}[H]
\centering
\fbox{\begin{minipage}{0.9\linewidth}\centering\vspace{6em}
\emph{[SCREENSHOT 2: Sidebar showing iris.csv selected as active,
services still green.]}\vspace{6em}\end{minipage}}
\caption{Post-upload sidebar.}
\label{fig:demo-2}
\end{figure}

\subsection*{Beat 3 (1:00--1:45) --- Ask for EDA}
In the chat, type:
\begin{quote}
\emph{``Run a full EDA on my data and summarise what stands out.''}
\end{quote}
Show the tool activity ribbon --- \tool{ingest\_dataset} then
\tool{run\_full\_eda} --- and the assistant's summary paragraphs.
Expand the tool-call cards to reveal the arguments and the JSON
result.

\begin{figure}[H]
\centering
\fbox{\begin{minipage}{0.9\linewidth}\centering\vspace{6em}
\emph{[SCREENSHOT 3: Chat with EDA tool cards expanded and a
correlation heatmap rendered.]}\vspace{6em}\end{minipage}}
\caption{EDA response with expanded tool cards.}
\label{fig:demo-3}
\end{figure}

\subsection*{Beat 4 (1:45--2:45) --- Train a model}
Type:
\begin{quote}
\emph{``Now train a model to predict Species and show me the
leaderboard.''}
\end{quote}
Show the live progress indicator counting up as
\tool{run\_parallel\_bake\_off} completes, then the assistant's
leaderboard summary. Emphasise that the champion is reported
verbatim from the tool result, not fabricated.

\subsection*{Beat 5 (2:45--3:15) --- Explain the model}
Type:
\begin{quote}
\emph{``Explain the champion --- which features matter most?''}
\end{quote}
Depending on the champion, the LLM will call either
\tool{calculate\_shap\_values} or
\tool{generate\_feature\_importance\_plot}. Show the resulting
importance plot inline.

\subsection*{Beat 6 (3:15--3:45) --- Generate a notebook + PDF}
Type:
\begin{quote}
\emph{``Generate a Jupyter notebook that reproduces this pipeline
and a PDF report of the results.''}
\end{quote}
Show \tool{generate\_jupyter\_notebook} and
\tool{compile\_pdf\_report} firing.

\subsection*{Beat 7 (3:45--4:30) --- Browse artefacts}
Switch to the \emph{Artifacts} tab. Show the filter chips, the
image previews for the plots, the parquet dataframe preview for the
\code{xtrain\_sample}, and the download buttons.

\begin{figure}[H]
\centering
\fbox{\begin{minipage}{0.9\linewidth}\centering\vspace{6em}
\emph{[SCREENSHOT 4: Artifacts tab showing at least six artefacts
across the plots, model, notebook, and report categories.]}\vspace{6em}\end{minipage}}
\caption{Artefacts tab after a full run.}
\label{fig:demo-4}
\end{figure}

\subsection*{Beat 8 (4:30--5:00) --- Open the downloads}
Cut to the OS file explorer showing \file{data/artifacts/plots/},
\file{data/artifacts/models/}, \file{data/artifacts/notebooks/},
and \file{data/artifacts/reports/}. Open the generated
\file{.ipynb} in Jupyter to prove it is executable, and the PDF in
a viewer to prove the leaderboard is present.

\section{Expected Terminal Output}
\label{sec:terminal-output}

For reference, the truncated log for a complete run resembles the
following (elided for brevity):

\begin{lstlisting}[style=alBash, numbers=none]
mcp.pool | Pool ready (in-process): 5/5 services online, 11 unified tools
mcp_data | Processing request of type CallToolRequest    # ingest
workspace | Artifact saved: kind=plot id=...             # correlation matrix
workspace | Artifact saved: kind=plot id=...             # distributions
workspace | Artifact saved: kind=plot id=...             # missing values
mcp_data | Processing request of type CallToolRequest    # bake-off
context.py Sending INFO to client: Detected problem type: classification
context.py Sending INFO to client: Training 4 models in parallel
context.py Sending INFO to client: Champion: LightGBM
workspace | Artifact saved: kind=model id=...
workspace | Artifact saved: kind=xtrain_sample id=...
workspace | Artifact saved: kind=plot id=...             # SHAP summary
workspace | Artifact saved: kind=training_notebook id=...
workspace | Artifact saved: kind=pdf_report id=...
\end{lstlisting}

\section{Post-Recording Editing Notes}
\label{sec:post-record}

Recommendations for the video edit:

\begin{itemize}[leftmargin=1.4em]
  \item Overlay the assistant's summary paragraphs at $1.5\times$
        playback while narrating; the raw generation speed varies
        with Groq load.
  \item Highlight the ``\code{artifact\_id}'' fields in the tool
        result to show that every artefact has a stable identifier.
  \item Show the size of \file{data/artifacts/\_index.json} growing
        after each artefact save --- this makes the registry
        tangible.
\end{itemize}
